% Generated draft: requires mathematical and source review.
\documentclass[11pt]{article}
\usepackage[margin=25mm]{geometry}
\usepackage{fontspec}
\setmainfont{Latin Modern Roman}
\usepackage{amsmath,amssymb,mathtools,unicode-math}
\setmathfont{Latin Modern Math}
\usepackage{enumitem,needspace,xurl,hyperref}
\setlength{\parindent}{0pt}
\setlength{\parskip}{5pt}
\setlength{\emergencystretch}{3em}
\newcommand{\N}{\mathbb{N}}
\newcommand{\Z}{\mathbb{Z}}
\newcommand{\Q}{\mathbb{Q}}
\newcommand{\R}{\mathbb{R}}
\newcommand{\C}{\mathbb{C}}
\newcommand{\F}{\mathbb{F}}
\newcommand{\A}{\mathbb{A}}
\newcommand{\PP}{\mathbb{P}}
\newcommand{\E}{\mathbb{E}}
\newcommand{\eps}{\varepsilon}
\newcommand{\dd}{\,\mathrm{d}}
\begin{document}
\section*{UnsolvedMath ID: 30002362}
\textbf{Source problem code:} catalogue entry
\subsection{Definitions and mathematical statement}
The supplied record concerns Lott's secondary K-group and uses the notation
\[
G(M):=KR^0(M),\qquad D(M):=KR_d^0(M),\qquad T(M):=KR_{\mathbb{R}/\mathbb{Z}}^{-1}(M).
\]
These are labels for the groups appearing in the source record, not independent definitions of them. The record does not specify the coefficient object \(R\), the class of admissible spaces \(M\), or the constructions defining these groups. It describes generators of \(G(M)\) by triples \((V,g_V,\eta)\), but does not supply their admissibility conditions or the defining relations. The operations \(\operatorname{cycl}\) and \(a\) below likewise retain their source notation; their domains and definitions are not supplied.

The extracted assertion is that there is a well-defined group homomorphism
\[
\Phi_M:G(M)\longrightarrow D(M),
\qquad D(M)\cong T(M),
\]
whose value on every admissible generator is
\[
\Phi_M\bigl([(V,g_V,\eta)]\bigr)
=\operatorname{cycl}(V,g_V)-a(\eta).
\]
Here brackets denote the class of the indicated generator in the secondary K-group. Well-definedness requires that the displayed expression belong to \(D(M)\) and respect every defining relation of \(G(M)\). The supplied record does not provide the hypotheses or relations needed to state those requirements explicitly, nor does it explain the asserted identification \(D(M)\cong T(M)\).

The supplied statement-verification field says that this extracted assertion is a construction and result proved in the cited report, rather than an unresolved question; it marks the statement as \texttt{unrecoverable} and supplies no clean replacement. Separately, the supplied dated literature assessment describes an existence part of Lott's conjecture as settled and a transfer-compatibility part as unproved. No transfer maps, compatibility equation, or hypotheses for that remaining part are supplied. Consequently, the displayed assertion preserves the extracted comparison-map claim, but the supplied evidence does not determine a precise remaining open problem.
\subsection{Short English statement}
The extracted claim sends each generator \((V,g_V,\eta)\) of Lott's secondary K-group to \(\operatorname{cycl}(V,g_V)-a(\eta)\) in the indicated differential algebraic K-theory group. The supplied verification calls this a proved result. Its separate dated assessment mentions unresolved transfer compatibility, but gives insufficient definitions and hypotheses to formulate that question precisely.
\subsection{Sources}
\begin{itemize}
\item Local UnsolvedMath registry, permanent ID \(30002362\), title \emph{Secondary K-Theory Comparison Map}, collection \emph{Oberwolfach Reports Open Problems}, exact source problem code \texttt{catalogue entry}. The record attributes the conjecture and secondary K-group to Lott; no individual proposer is supplied. Its proposed-year field is \(2013\).
\item Supplied report citation: \emph{Algebraic K-theory and Motivic Cohomology} (2014), Oberwolfach Reports, pp.\ 1861--1913, DOI \texttt{10.4171/owr/2013/32}. \url{https://doi.org/10.4171/owr/2013/32}. The background describes these reports as documenting mathematical workshops and their open-problem discussions.
\item Supplied literature source, labeled ``Oberwolfach report on Lott's conjecture'': \url{https://ems.press/content/serial-article-files/46461?nt=1}.
\item Supplied literature source, labeled ``Bunke and Gepner, differential K-theory transfer'': \url{https://arxiv.org/abs/1306.0247}.
\item Literature assessment dated exactly \texttt{2026-08-21}: the supplied curation labels the problem partially solved and states that the differential algebraic K-theory theorem settles the existence part of Lott's conjecture, while transfer compatibility remains unproved. This is supplied dated curation derived from the OWR record, not an independently verified literature assessment; its time-sensitive status and cited claims require independent verification.
\item The supplied statement-verification field reports a check against the official Oberwolfach report and classifies the extracted sentence as a proved construction/result. It supplies no separate clean statement and retains the original record for provenance. This verification report is reproduced as supplied metadata, not as a check performed for this draft.
\end{itemize}
\end{document}
